% TEMPLATE-MILC-v3.tex - plantilla maestra para todas las guias MILC v3
% Ver scripts/generadores/build-guia-milc.py para la lista completa de
% claves esperadas. El script valida que no queden placeholders sin
% reemplazar antes de escribir el archivo final.

\documentclass[11pt,letterpaper]{article}
\usepackage[margin=1.75cm]{geometry}
\usepackage[table]{xcolor}
\usepackage{fontspec}
\usepackage[spanish]{babel}
\usepackage{tikz}
\usepackage[most]{tcolorbox}
\usepackage{tabularx}
\usepackage{array}
\usepackage{fancyhdr}
\usepackage{titlesec}
\usepackage{ragged2e}
\usepackage{enumitem}
\usepackage{microtype}

\setmainfont{Helvetica Neue}
\setsansfont{Helvetica Neue}
\setlength{\parindent}{0pt}
\setlength{\parskip}{6pt}
\hyphenpenalty=200
\exhyphenpenalty=200
\pretolerance=400
\tolerance=2000
\setlength{\emergencystretch}{1em}
\renewcommand{\arraystretch}{1.25}
\setlist{leftmargin=5mm,itemsep=1mm,topsep=1mm}
\newcolumntype{Y}{>{\RaggedRight\arraybackslash}X}

\definecolor{milcMagenta}{HTML}{E6007E}
\definecolor{milcVino}{HTML}{5A0038}
\definecolor{milcTurquesa}{HTML}{0093A5}
\definecolor{milcVerde}{HTML}{58A923}
\definecolor{milcMostaza}{HTML}{E5B400}
\definecolor{milcOcre}{HTML}{B86600}
\definecolor{milcNegro}{HTML}{241816}
\definecolor{milcGris}{HTML}{F7F7F4}
\definecolor{milcLinea}{HTML}{E8E2DD}
\definecolor{milcGradePrimary}{HTML}{ED7446}
\definecolor{milcGradeDark}{HTML}{B63E16}
\definecolor{milcGradeSoft}{HTML}{FFF0E8}
\definecolor{milcGradeAccent}{HTML}{CFE7FF}
\definecolor{milcGradeText}{HTML}{FFFFFF}
\color{milcNegro}

\pagestyle{fancy}
\fancyhf{}
\lhead{\small Tecnología e Informática · Grado 8}
\rhead{\small Guía 8 · MILC v3}
\cfoot{\small\thepage}
\renewcommand{\headrulewidth}{0pt}

\newtcolorbox{titlebox}[1]{
  enhanced,colback=#1,colframe=#1,arc=7mm,boxrule=0pt,
  left=7mm,right=7mm,top=3mm,bottom=3mm,
  before skip=8pt,after skip=8pt,
  fontupper=\sffamily\bfseries\Large\color{white}
}

\newtcolorbox{softbox}[3]{
  enhanced,colback=#2,colframe=#1,borderline west={4pt}{0pt}{#1},
  arc=2mm,boxrule=.4pt,left=4mm,right=4mm,top=3mm,bottom=3mm,
  before skip=6pt,after skip=8pt,title={#3},coltitle=white,
  colbacktitle=#1,fonttitle=\sffamily\bfseries,fontupper=\RaggedRight,
  attach boxed title to top left={xshift=3mm,yshift=-2mm},
  boxed title style={arc=2mm, boxrule=0pt}
}

\newtcolorbox{drawbox}[2]{
  enhanced,colback=white,colframe=#1,arc=4mm,boxrule=1pt,
  left=5mm,right=5mm,top=4mm,bottom=4mm,before skip=8pt,after skip=8pt,
  title={#2},coltitle=white,colbacktitle=#1,fonttitle=\sffamily\bfseries,
  fontupper=\RaggedRight,
  attach boxed title to top left={xshift=4mm,yshift=-2mm},
  boxed title style={arc=3mm, boxrule=0pt}
}

\newtcolorbox{infoband}[1]{
  enhanced,colback=#1!8,colframe=#1!50,arc=2mm,boxrule=.5pt,
  left=4mm,right=4mm,top=2mm,bottom=2mm,before skip=4pt,after skip=4pt,
  fontupper=\small\RaggedRight
}

\newcommand{\linead}[1]{\noindent\color{milcLinea}\rule{#1}{0.5pt}\color{milcNegro}}
\newcommand{\respuesta}[1]{\par\vspace{2mm}\foreach \n in {1,...,#1}{\noindent\color{milcLinea}\rule{\linewidth}{0.5pt}\par\vspace{5mm}}\color{milcNegro}}
\newcommand{\checkbox}{\textcolor{milcGradeDark}{\fbox{\phantom{x}}}\hspace{5pt}}

\newcommand{\phasecard}[4]{
\begin{tcolorbox}[enhanced,colback=#3,colframe=#2,arc=3mm,boxrule=.5pt,height=3.05cm,valign=center,halign=center,left=1.5mm,right=1.5mm,top=2mm,bottom=2mm]
{\sffamily\bfseries\fontsize{22}{24}\selectfont\textcolor{#2}{#1}}\par
{\sffamily\bfseries\small #4}
\end{tcolorbox}
}

\begin{document}
\thispagestyle{empty}
\begin{tikzpicture}[remember picture,overlay]
  \fill[white] (current page.south west) rectangle (current page.north east);
  \path[fill=milcGradePrimary,rounded corners=16mm]
    ([xshift=.55cm,yshift=-.55cm]current page.north west) rectangle
    ([xshift=-.55cm,yshift=3.65cm]current page.south east);

  \node[anchor=north west,text=milcGradeText] at ([xshift=2.0cm,yshift=-1.85cm]current page.north west)
    {\sffamily\bfseries\fontsize{14}{16}\selectfont GRADO};
  \node[anchor=north west,text=milcGradeText,opacity=.82] at ([xshift=2.0cm,yshift=-2.75cm]current page.north west)
    {\sffamily\bfseries\fontsize{9}{11}\selectfont SERIE GUÍAS · TECNOLOGÍA E INFORMÁTICA};

  \begin{scope}[shift={([xshift=-3.05cm,yshift=-2.55cm]current page.north east)}]
    \fill[white,opacity=.80] (-.62,-.52) rectangle (.62,.52);
    \draw[milcGradeText,line width=1pt] (-.62,-.52) rectangle (.62,.52);
    \fill[milcGradeDark] (-.42,-.38) rectangle (-.22,.06);
    \fill[milcGradeAccent] (-.10,-.38) rectangle (.10,.24);
    \fill[milcGradePrimary!60!black] (.22,-.38) rectangle (.42,.38);
  \end{scope}

  \node[anchor=west,text=milcGradeText] at ([xshift=1.9cm,yshift=-12.85cm]current page.north west)
    {\sffamily\bfseries\fontsize{124}{124}\selectfont 8};
  \draw[milcGradeText,line width=1.7pt]
    ([xshift=4.52cm,yshift=-11.68cm]current page.north west) circle (.13cm);

  \node[anchor=west,text=milcGradeText,text width=8.7cm,align=left] at ([xshift=10.35cm,yshift=-11.6cm]current page.north west)
    {
     {\sffamily\bfseries\fontsize{19}{22}\selectfont Formato condicional y validación de datos — presentación profesional de información}\\[4mm]
     {\sffamily\bfseries\fontsize{23}{26}\selectfont Guía 8-8-TIC}\\[2mm]
     {\sffamily\fontsize{13}{16}\selectfont Periodo Primero}\\[5mm]
     {\sffamily\bfseries\fontsize{11}{13}\selectfont Institución Educativa Sor María Juliana}\\[1mm]
     {\sffamily\fontsize{10}{12}\selectfont Autor: Dr. Álvaro Cárdenas Orozco}
    };

  \path[fill=milcGradeSoft,rounded corners=7mm]
    ([xshift=1.35cm,yshift=.85cm]current page.south west) rectangle
    ([xshift=-1.35cm,yshift=3.15cm]current page.south east);
  \draw[milcGradeDark,line width=.65pt,rounded corners=7mm]
    ([xshift=1.35cm,yshift=.85cm]current page.south west) rectangle
    ([xshift=-1.35cm,yshift=3.15cm]current page.south east);
  \node[anchor=center,text=milcNegro,text width=17.0cm,align=center] at ([yshift=2.0cm]current page.south)
    {\sffamily\fontsize{10}{12}\selectfont Metodología MILC · Apertura ancestral · Pensamiento computacional · Triángulo Dussel-Estoicismo-Floridi\\
    \textcolor{milcGradeDark}{\bfseries Producto:} Hoja con lista de asistencia validada (P/A obligatorio) + formato condicional rojo/verde + protección de fila de totales.};
\end{tikzpicture}
\mbox{}
\newpage

\section*{Apertura: Formato condicional y validación de datos — presentación profesional de información}

\begin{softbox}{milcGradeDark}{milcGradeSoft}{Saber ancestral}
El \textbf{semáforo del salón comunal}: tres colores, tres reglas. Verde = pasa, amarillo = espera, rojo = no entres. Cada color tiene regla y todos en la comunidad la entienden al instante. El portero del salón también: si no traes invitación, no pasas. Reglas visuales y reglas de entrada — disciplina comunitaria.
\end{softbox}

\begin{softbox}{milcTurquesa}{milcGris}{Saber contemporáneo}
El \textbf{formato condicional} de Excel cambia el color de la celda según una regla: si la nota es menor a 3.0, se pinta roja; si es mayor a 4.5, verde. La \textbf{validación de datos} restringe qué se puede escribir: solo letras P o A, solo fechas, solo números entre 0 y 5.
\end{softbox}

\begin{softbox}{milcMostaza}{milcGris}{Pregunta-puente}
\textit{¿Qué del semáforo del salón comunal — disciplina visual heredada — podemos llevar a la hoja de Excel para que ayude a leer y a no equivocarse?}
\end{softbox}

\begin{softbox}{milcVerde}{milcGris}{Saber y saber hacer}
Al terminar podrás \textbf{aplicar formato condicional} con reglas claras (mayor que, menor que, contiene), \textbf{configurar validación de datos} con listas y rangos, \textbf{proteger celdas críticas} y \textbf{diseñar hojas autovalidadas} que orientan al usuario.
\end{softbox}

\small
\begin{tabularx}{\textwidth}{>{\bfseries}p{3.0cm}Y}
\rowcolor{milcGradePrimary!12}Autor & Dr. Álvaro Cárdenas Orozco\\
DBA & Diseña hojas con formato condicional y validación de datos para presentación profesional y prevención de errores.\\
Referentes & MEN Tecnología · ISO 9001 (calidad) · Modelo MILC · Dussel (justicia, no poder) · Epicteto (lo que depende de mí)\\
Producto final & Hoja con lista de asistencia validada (P/A obligatorio) + formato condicional rojo/verde + protección de fila de totales.\\
\end{tabularx}
\normalsize

\begin{titlebox}{milcGradeDark}
Ruta MILC: así vamos a aprender
\end{titlebox}
\begin{tabularx}{\textwidth}{>{\centering\arraybackslash}X>{\centering\arraybackslash}X>{\centering\arraybackslash}X>{\centering\arraybackslash}X}
\phasecard{1}{milcVerde}{milcGris}{Escuta\\[-1mm]\small Observo, escucho y pregunto.} &
\phasecard{2}{milcTurquesa}{milcGris}{Sistematización\\[-1mm]\small Pensamiento computacional.} &
\phasecard{3}{milcMagenta}{milcGris}{Praxis\\[-1mm]\small Construyo y aplico.} &
\phasecard{4}{milcVino}{milcGris}{Evaluación\\[-1mm]\small Reflexión liberadora.}
\end{tabularx}


\begin{titlebox}{milcVerde}
Fase 1 · Escuta
\end{titlebox}

\begin{softbox}{milcVerde}{milcGris}{Escena para escuchar}
Daniela hace la lista de asistencia del curso en Excel. Le piden que las faltas (A) se vean en rojo y los presentes (P) en verde. Y que ningún estudiante pueda escribir un valor distinto a P o A. Y que la fila de totales (con fórmulas) NO se pueda modificar por error. ¿Cómo configura todo esto en Excel?
\end{softbox}

\checkbox Recuerdo una hoja de Excel donde alguien escribió mal y dañó los cálculos.\\
\checkbox Identifico cuándo conviene formato condicional y cuándo validación de datos.\\
\checkbox Distingo entre proteger una celda y ocultar la fórmula.

\textbf{Una vez una hoja se dañó porque alguien escribió:}\\[1mm]
\linead{\linewidth}\\[1mm]
\linead{\linewidth}

\textbf{Lo que se podría haber configurado para evitar el error es:}\\[1mm]
\linead{\linewidth}\\[1mm]
\linead{\linewidth}

\textbf{El semáforo digital (formato condicional) lo aplicaría así en mi vida:}\\[1mm]
\linead{\linewidth}\\[1mm]
\linead{\linewidth}

\begin{infoband}{milcVerde}
\textbf{Saber más.} El formato condicional y la validación son las dos herramientas que separan una hoja amateur de una hoja profesional. Empresas serias las usan en TODAS sus plantillas: nómina, inventario, presupuesto. Sin ellas, los errores humanos contaminan los datos.
\end{infoband}

\begin{titlebox}{milcTurquesa}
Fase 2 · Sistematización con pensamiento computacional
\end{titlebox}

\begin{softbox}{milcTurquesa}{milcGris}{Cuatro pilares aplicados al tema}
El pensamiento computacional descompone la hoja en sus reglas visuales y de entrada, reconoce patrones de error humano, abstrae la regla \textit{prevenir es mejor que corregir}, y diseña un algoritmo de protección.
\end{softbox}

\small
\begin{tabularx}{\textwidth}{>{\bfseries\RaggedRight\arraybackslash}p{3.6cm}Y}
\rowcolor{milcTurquesa!90}\color{white}Pilar & \color{white}\bfseries Aplicación al tema\\
1. Descomposición & Descomponer la hoja: ¿qué celdas necesitan resaltado visual? ¿qué celdas necesitan limitar entrada? ¿qué celdas hay que proteger?\\
2. Reconocimiento de patrones & Patrones: notas bajas $\rightarrow$ rojo. Aciertos $\rightarrow$ verde. Lista de opciones cerrada (P/A) $\rightarrow$ validación con lista. Fórmulas críticas $\rightarrow$ proteger.\\
3. Abstracción & Abstracción: la hoja debe enseñarle al usuario cómo usarla — sin necesidad de manual.\\
4. Algoritmo conceptual & Algoritmo: \textbf{1)} identificar celdas a resaltar · \textbf{2)} aplicar formato condicional · \textbf{3)} identificar celdas a restringir · \textbf{4)} configurar validación · \textbf{5)} proteger fórmulas · \textbf{6)} probar con datos de prueba.\\
\end{tabularx}
\normalsize

\begin{drawbox}{milcTurquesa}{Anatomía de una hoja autovalidada}
\textbf{Formato condicional:} reglas visuales (color por valor, ícono, barra de datos). \\
\textbf{Validación de datos:} restricciones de entrada (lista, rango numérico, longitud). \\
\textbf{Mensaje de entrada:} pista al usuario antes de escribir. \\
\textbf{Mensaje de error:} explicación cuando se intenta valor inválido. \\
\textbf{Protección de hoja:} bloquear celdas críticas (con o sin contraseña). \\
\textbf{Verificación:} probar con un valor válido y uno inválido para confirmar.
\end{drawbox}

\begin{softbox}{milcMostaza}{milcGris}{Errores comunes y su corrección}
\begin{itemize}
\item Aplicar formato condicional sin verificar la regla $\rightarrow$ los colores no reflejan lo esperado.
\item Validación con lista pero sin mensaje de error $\rightarrow$ usuario no entiende por qué Excel rechaza.
\item Proteger TODA la hoja $\rightarrow$ nadie puede usarla; protege SOLO las celdas críticas.
\item Olvidar desproteger antes de modificar fórmulas $\rightarrow$ frustrante para el dueño.
\item Confundir formato con valor $\rightarrow$ el formato es máscara visual, el valor es lo real.
\end{itemize}
\end{softbox}

\textbf{Para que las notas menores a 3.0 se vean en rojo automáticamente, configuraría:}\\[1mm]
\linead{\linewidth}\\[1mm]
\linead{\linewidth}

\textbf{Para que solo se acepte P o A en una columna, configuraría:}\\[1mm]
\linead{\linewidth}\\[1mm]
\linead{\linewidth}

\begin{titlebox}{milcMagenta}
Fase 3 · Praxis con pensamiento computacional
\end{titlebox}

\begin{softbox}{milcMagenta}{milcGris}{Construyendo el producto con los cuatro pilares}
Construir una hoja autovalidada es un sistema de prevención. Decomponemos las celdas por riesgo, reconocemos patrones de error, abstraemos la regla \textit{ayudar al usuario antes que corregirlo}, y seguimos un algoritmo paso a paso.
\end{softbox}

\small
\begin{tabularx}{\textwidth}{>{\bfseries\RaggedRight\arraybackslash}p{3.6cm}Y}
\rowcolor{milcMagenta!90}\color{white}Pilar & \color{white}\bfseries Aplicación al producto\\
Descomposición & Hoja: 1) cabecera de la lista · 2) columna de nombres protegida · 3) columna de asistencia validada con P/A · 4) formato condicional rojo/verde · 5) fila de totales con fórmulas protegidas.\\
Patrones & Patrones de hojas profesionales: validación + formato + protección actuando juntos para que el usuario no pueda dañar la hoja.\\
Abstracción & Función esencial: registrar asistencia rápido y sin errores, con visual claro.\\
Algoritmo de armado & Pasos: crear lista $\rightarrow$ aplicar validación P/A $\rightarrow$ formato condicional verde/rojo $\rightarrow$ fórmulas de totales $\rightarrow$ protección selectiva $\rightarrow$ prueba con datos.\\
\end{tabularx}
\normalsize

\begin{drawbox}{milcMagenta}{Checklist de la lista de asistencia profesional}
\checkbox Lista con nombres de 10+ estudiantes.\\
\checkbox Validación de datos: solo P o A en columna de asistencia.\\
\checkbox Formato condicional: P en verde, A en rojo.\\
\checkbox Mensaje de entrada (\textit{``Escribe P o A''}).\\
\checkbox Mensaje de error si se intenta otro valor.\\
\checkbox Fila de totales con fórmulas =CONTAR.SI.\\
\checkbox Protección selectiva: bloquear solo celdas de totales.
\end{drawbox}

\begin{softbox}{milcMagenta}{milcGris}{Plantilla de trabajo}
\textbf{Lista:} columna A nombres, columna B asistencia. \\
\textbf{Validación columna B:} Datos $\rightarrow$ Validación $\rightarrow$ Lista $\rightarrow$ \textit{P,A}. \\
\textbf{Formato condicional B:} Inicio $\rightarrow$ Formato condicional $\rightarrow$ Reglas: \textit{=B2=``A''} $\rightarrow$ rojo; \textit{=B2=``P''} $\rightarrow$ verde. \\
\textbf{Totales:} =CONTAR.SI(B2:B11,``P'') y =CONTAR.SI(B2:B11,``A''). \\
\textbf{Protección:} Revisar $\rightarrow$ Proteger hoja, dejando libre solo la columna B.
\end{softbox}

\textbf{Mi hoja de asistencia con validación y formato condicional:}\\[1mm]
\linead{\linewidth}\\[1mm]
\linead{\linewidth}\\[1mm]
\linead{\linewidth}

\begin{titlebox}{milcMostaza}
Producto de la sesión
\end{titlebox}
\textbf{Lista de asistencia validada + formato condicional + protección}.
\begin{softbox}{milcMostaza}{milcGris}{Criterios de calidad}
\begin{itemize}
\item Lista con 10+ estudiantes reales o simulados.
\item Validación de datos funcionando (rechaza valores distintos a P/A).
\item Formato condicional verde/rojo funcionando.
\item Mensajes de entrada y error.
\item Totales con fórmulas =CONTAR.SI.
\item Protección selectiva (no toda la hoja).
\end{itemize}
\end{softbox}

\begin{titlebox}{milcVino}
Fase 4 · Evaluación liberadora — cinco dimensiones
\end{titlebox}

\begin{softbox}{milcVino}{milcGris}{Sentido de la evaluación}
Evaluar no es solo revisar una nota. Es reconocer qué aprendí, qué cambió en mí, qué emociones manejé, cómo me relaciono con la ciudad y la comunidad, y qué le dejaré a quien viene detrás.
\end{softbox}

\textbf{Mi reflexión liberadora en cinco dimensiones.} Completa cada frase iniciada:

\small
\begin{tabularx}{\textwidth}{>{\bfseries\RaggedRight\arraybackslash}p{3.4cm}Y>{\RaggedRight\arraybackslash}p{4.6cm}}
\rowcolor{milcVino}\color{white}Dimensión & \color{white}\bfseries Mi respuesta (frame starter) & \color{white}\bfseries Algo que descubrí\\
1. Personal & Lo que más cambió en mí fue \linead{6.5cm} & \linead{4.2cm}\\
2. Emocional & La emoción más fuerte fue \linead{2.5cm} y la regulé \linead{3cm} & \linead{4.2cm}\\
3. Ciudadana & En democracia esto importa porque \linead{6cm} & \linead{4.2cm}\\
4. Local (Cartago) & En mi barrio sería útil para \linead{6.5cm} & \linead{4.2cm}\\
5. Intergeneracional & A mi abuela/abuelo le diría \linead{6.5cm} & \linead{4.2cm}\\
\end{tabularx}
\normalsize

\begin{infoband}{milcVino}
\textbf{Para la próxima sustentación.} Vuelve a esta tabla en seis meses. Verás que las respuestas cambian — no porque lo aprendido cambie, sino porque tú cambias. La evaluación liberadora es una conversación contigo a través del tiempo.
\end{infoband}


\begin{titlebox}{milcGradeDark}
Triángulo de pensamiento · cierre filosófico
\end{titlebox}

\begin{softbox}{milcVino}{milcGris}{Dussel · lente del nosotros}
\textit{``Sin reconocimiento del Otro no hay justicia, sólo poder.''}

La validación de datos respeta a quien va a llenar la hoja: le dice qué se espera, no le castiga después por no saber. Es ética del diseño: hacer la herramienta para el usuario, no contra él.

\textbf{¿Tu hoja ayuda al usuario o lo culpa por no saber qué escribir?}
\end{softbox}
\linead{\linewidth}\\[1mm]
\linead{\linewidth}

\begin{softbox}{milcMostaza}{milcGris}{Estoicismo · Epicteto · cuidado interior}
\textit{``Algunas cosas dependen de nosotros y otras no — distinguir esto es la libertad.''}

El formato condicional automatiza lo que controlas: el color según el valor. La validación restringe lo que el usuario puede escribir. Lo que cada uno hace dentro de esos límites — escribir P o A — sigue siendo decisión suya.

\textbf{¿Qué de tu hoja automatizaste tú y qué dejaste libre para el usuario? ¿Es una decisión consciente?}
\end{softbox}
\linead{\linewidth}\\[1mm]
\linead{\linewidth}

\begin{softbox}{milcTurquesa}{milcGris}{Floridi · lente de la infoesfera}
\textit{``La información es un bien común; merece el cuidado del agua o del aire.''}

La validación protege la calidad del bien común: una tabla limpia. Sin validación, cualquier usuario puede contaminar los datos por descuido. La hoja autovalidada cuida la información de todos.

\textbf{Si tu hoja la usaran 30 personas, ¿la validación protegería la integridad de los datos?}
\end{softbox}
\linead{\linewidth}\\[1mm]
\linead{\linewidth}

\begin{titlebox}{milcVino}
Compromiso final
\end{titlebox}

\small
\begin{tabularx}{\textwidth}{>{\RaggedRight\arraybackslash}p{3.0cm}YYY}
\rowcolor{milcVino}\color{white}\bfseries Criterio & \color{white}\bfseries Lo logré & \color{white}\bfseries En proceso & \color{white}\bfseries Necesito apoyo\\
Comprensión del tema & Explico ideas centrales con mis palabras. & Reconozco ideas pero falta precisión. & Necesito repasar conceptos base.\\
Pensamiento computacional & Apliqué los 4 pilares al diseñar mi producto. & Apliqué algunos pilares. & Necesito releer la fase 2.\\
Aplicación y producto & Mi producto cumple los criterios y se puede revisar. & Producto funcional con ajustes pendientes. & Aún en construcción.\\
Reflexión 5 dimensiones & Respondí cada dimensión con honestidad. & Respondí algunas con detalle. & Mis respuestas son muy generales.\\
Triángulo de pensamiento & Conecté las 3 voces con mi trabajo real. & Conecté al menos una. & No relaciono las voces con mi proceso.\\
\end{tabularx}
\normalsize

\textbf{Hoy entendí que:}\\[1mm]
\linead{\linewidth}\\[1mm]
\linead{\linewidth}

\textbf{Mi compromiso para la próxima sesión es:}\\[1mm]
\linead{\linewidth}\\[1mm]
\linead{\linewidth}

\begin{softbox}{milcMostaza}{milcGris}{Cita de despedida · Marco Aurelio}
\textit{``El obstáculo en el camino se vuelve el camino. Lo que estorba al camino se vuelve camino.''}\\[1mm]
Cada error de hoy, bien observado, se convierte en pieza del camino que pisarás mañana.
\end{softbox}

\small
\begin{tabularx}{\textwidth}{>{\bfseries}p{3.6cm}Y}
\rowcolor{milcGradeDark!12}Nombre completo: & \linead{10cm}\\
Fecha: & \linead{4cm} \quad Curso/grupo: \linead{4cm}\\
Mi firma: & \linead{9cm}\\
\end{tabularx}
\normalsize

\begin{infoband}{milcGradeDark}
\textbf{Próxima sesión.} Aplicaremos todo lo aprendido a un mini-estudio de datos del entorno: pregunta, recolección, análisis y conclusión prudente.
\end{infoband}

\end{document}
